% GENERATED FILE. Edit canonical lessons, then run npm run generate:books.
% canonical-source-hash: fnv1a64:fdf268673ed36bb0
% canonical-lessons: TE-C107-korika, TE-C107-asa, TE-C107-kala, TE-C107-uddesam, TE-C107-pranalika

\chapter{Wish, Hope, Dream, Intention, Plan}
\label{ch:te-wish-hope-dream-intention-plan}
\hlchaptermodality{107}

\begin{chapteropening}
\textbf{By the end of this chapter:} \emph{I can say a wish, hope, a dream, an intention and a plan.}

Complete the last lesson of chapter 107: I can say a wish, hope, a dream, an intention and a plan.
\end{chapteropening}

\section[\texorpdfstring{\te{కోరిక} (kōrika)}{కోరిక (kōrika)}]{\te{కోరిక} (kōrika) — a wish}

\label{lesson:TE-C107-korika}



\begin{quote}
Before the new one: say the Telugu for an owner, then the Telugu for strength.
\end{quote}

\subsection*{You'll want to know: \te{కోరిక}}
\textbf{\te{కోరిక}} — \emph{kōrika} — \textquotedblleft{}a wish\textquotedblright{}.

\begin{grammarlens}[title={What you've built}]
One more word for this chapter.
\end{grammarlens}

\subsection*{Your turn}
\begin{itemize}
\raggedright
  \item \emph{Say it:} \emph{kōrika}
  \item \emph{Say it:} \emph{kōrika}, once more
  \item \emph{Say it:} say \emph{kōrika}
  \item \emph{Recall:} say the Telugu for an owner, then the Telugu for strength, then say \emph{kōrika} again
\end{itemize}

\begin{tcolorbox}[breakable,colback=teal!4,colframe=teal!35!black,title={Before you move on}]
What does \te{కోరిక} mean? (\textquotedblleft{}A wish\textquotedblright{}.) Say it once more.
\end{tcolorbox}
\section[\texorpdfstring{\te{ఆశ} (āśa)}{ఆశ (āśa)}]{\te{ఆశ} (āśa) — hope}

\label{lesson:TE-C107-asa}



\begin{quote}
Before the new one: say the Telugu for strength, then the Telugu for a wish.
\end{quote}

\subsection*{You'll want to know: \te{ఆశ}}
\textbf{\te{ఆశ}} — \emph{āśa} — \textquotedblleft{}hope\textquotedblright{}.

\begin{grammarlens}[title={What you've built}]
One more word for this chapter.
\end{grammarlens}

\subsection*{Your turn}
\begin{itemize}
\raggedright
  \item \emph{Say it:} \emph{āśa}
  \item \emph{Say it:} \emph{āśa}, once more
  \item \emph{Say it:} say \emph{āśa}
  \item \emph{Recall:} say the Telugu for strength, then the Telugu for a wish, then say \emph{āśa} again
\end{itemize}

\begin{tcolorbox}[breakable,colback=teal!4,colframe=teal!35!black,title={Before you move on}]
What does \te{ఆశ} mean? (\textquotedblleft{}Hope\textquotedblright{}.) Say it once more.
\end{tcolorbox}
\section[\texorpdfstring{\te{కల} (kala)}{కల (kala)}]{\te{కల} (kala) — a dream}

\label{lesson:TE-C107-kala}



\begin{quote}
Before the new one: say the Telugu for a wish, then the Telugu for hope.
\end{quote}

\subsection*{You'll want to know: \te{కల}}
\textbf{\te{కల}} — \emph{kala} — \textquotedblleft{}a dream\textquotedblright{}.

\begin{grammarlens}[title={What you've built}]
One more word for this chapter.
\end{grammarlens}

\subsection*{Your turn}
\begin{itemize}
\raggedright
  \item \emph{Say it:} \emph{kala}
  \item \emph{Say it:} \emph{kala}, once more
  \item \emph{Say it:} say \emph{kala}
  \item \emph{Recall:} say the Telugu for a wish, then the Telugu for hope, then say \emph{kala} again
\end{itemize}

\begin{tcolorbox}[breakable,colback=teal!4,colframe=teal!35!black,title={Before you move on}]
What does \te{కల} mean? (\textquotedblleft{}A dream\textquotedblright{}.) Say it once more.
\end{tcolorbox}
\section[\texorpdfstring{\te{ఉద్దేశం} (uddēśaṁ)}{ఉద్దేశం (uddēśaṁ)}]{\te{ఉద్దేశం} (uddēśaṁ) — an intention}

\label{lesson:TE-C107-uddesam}



\begin{quote}
Before the new one: say the Telugu for hope, then the Telugu for a dream.
\end{quote}

\subsection*{You'll want to know: \te{ఉద్దేశం}}
\textbf{\te{ఉద్దేశం}} — \emph{uddēśaṁ} — \textquotedblleft{}an intention\textquotedblright{}.

\begin{grammarlens}[title={What you've built}]
One more word for this chapter.
\end{grammarlens}

\subsection*{Your turn}
\begin{itemize}
\raggedright
  \item \emph{Say it:} \emph{uddēśaṁ}
  \item \emph{Say it:} \emph{uddēśaṁ}, once more
  \item \emph{Say it:} say \emph{uddēśaṁ}
  \item \emph{Recall:} say the Telugu for hope, then the Telugu for a dream, then say \emph{uddēśaṁ} again
\end{itemize}

\begin{tcolorbox}[breakable,colback=teal!4,colframe=teal!35!black,title={Before you move on}]
What does \te{ఉద్దేశం} mean? (\textquotedblleft{}An intention\textquotedblright{}.) Say it once more.
\end{tcolorbox}
\section[\texorpdfstring{\te{ప్రణాళిక} (praṇālika)}{ప్రణాళిక (praṇālika)}]{\te{ప్రణాళిక} (praṇālika) — a plan}

\label{lesson:TE-C107-pranalika}



\begin{quote}
Before the new one: say the Telugu for a dream, then the Telugu for an intention.
\end{quote}

\subsection*{You'll want to know: \te{ప్రణాళిక}}
\textbf{\te{ప్రణాళిక}} — \emph{praṇālika} — \textquotedblleft{}a plan\textquotedblright{}.

\begin{grammarlens}[title={What you've built}]
One more word for this chapter.
\end{grammarlens}

\subsection*{Your turn}
\begin{itemize}
\raggedright
  \item \emph{Say it:} \emph{praṇālika}
  \item \emph{Say it:} \emph{praṇālika}, once more
  \item \emph{Say it:} say \emph{praṇālika}
  \item \emph{Recall:} say the Telugu for a dream, then the Telugu for an intention, then say \emph{praṇālika} again
\end{itemize}

\begin{tcolorbox}[breakable,colback=teal!4,colframe=teal!35!black,title={Before you move on}]
What does \te{ప్రణాళిక} mean? (\textquotedblleft{}A plan\textquotedblright{}.) Say it once more.
\end{tcolorbox}
